% ==============================================================================
% Paper 85 (DE): Lehmers Problem: Mahler-Maß und Algebraische Zahlentheorie
% Autor: Michael Fuhrmann
% Projekt: Specialist Mathematics Project, Build 167
% Datum: 2026-03-12
% ==============================================================================

\documentclass[12pt,a4paper]{amsart}

\usepackage{amsmath,amssymb,amsthm,mathtools,enumitem,booktabs,array,hyperref}
\usepackage[utf8]{inputenc}
\usepackage[T1]{fontenc}
\usepackage[ngerman]{babel}
\usepackage{geometry}
\geometry{margin=2.5cm}

% --------------------------------------------------------------------------
% Theorem-Umgebungen (Deutsch)
% --------------------------------------------------------------------------
\newtheorem{theorem}{Theorem}[section]
\newtheorem{lemma}[theorem]{Lemma}
\newtheorem{corollary}[theorem]{Korollar}
\newtheorem{proposition}[theorem]{Proposition}
\newtheorem{satz}[theorem]{Satz}
\newtheorem{korollar}[theorem]{Korollar}
\theoremstyle{definition}
\newtheorem{definition}[theorem]{Definition}
\newtheorem{definition_de}[theorem]{Definition}
\newtheorem{example}[theorem]{Beispiel}
\theoremstyle{remark}
\newtheorem{remark}[theorem]{Bemerkung}
\newtheorem{bemerkung}[theorem]{Bemerkung}
\newtheorem{conjecture}[theorem]{Vermutung}
\newtheorem{vermutung}[theorem]{Vermutung}

% --------------------------------------------------------------------------
% Makros
% --------------------------------------------------------------------------
\newcommand{\Z}{\mathbb{Z}}
\newcommand{\Q}{\mathbb{Q}}
\newcommand{\R}{\mathbb{R}}
\newcommand{\C}{\mathbb{C}}
\newcommand{\T}{\mathbb{T}}
\newcommand{\N}{\mathbb{N}}
\newcommand{\M}{\mathrm{M}}
\newcommand{\m}{\mathrm{m}}
\DeclareMathOperator{\deg}{deg}
\DeclareMathOperator{\Res}{Res}

\hypersetup{
  colorlinks=true,
  linkcolor=blue,
  citecolor=blue,
  urlcolor=blue,
  pdftitle={Lehmers Problem: Mahler-Maß und Algebraische Zahlentheorie},
  pdfauthor={Michael Fuhrmann}
}

% ==============================================================================
\begin{document}
% ==============================================================================

\title{Lehmers Problem:\\
  Mahler-Maß und Algebraische Zahlentheorie}

\author{Michael Fuhrmann}
\address{Specialist Mathematics Project, Build 167}
\email{specialist-maths@project.local}
\date{2026-03-12}

\begin{abstract}
Lehmers Problem, 1933 von D.\,H.\ Lehmer aufgeworfen, fragt, ob das Mahler-Maß
$\M(\alpha)$ einer algebraischen ganzen Zahl $\alpha$, die keine Einheitswurzel ist,
beliebig nahe an $1$ liegen kann. Das Mahler-Maß eines Polynoms $P(x)$ ist
\[
  \M(P) = \exp\!\left(\int_0^1 \log|P(e^{2\pi i t})|\,dt\right)
        = |a_d| \prod_{j=1}^d \max(1, |\alpha_j|),
\]
wobei $a_d$ der Leitkoeffizient und $\alpha_1,\ldots,\alpha_d$ die Nullstellen sind.
Lehmers Polynom $L(x) = x^{10}+x^9-x^7-x^6-x^5-x^4-x^3+x+1$ mit dem Mahler-Maß
$\M(L) \approx 1{,}17628\ldots$ (Salem-Zahl $\lambda_{10}$) ist trotz 90 Jahren
intensiver Suche das kleinste bekannte Mahler-Maß $> 1$ für algebraische ganze Zahlen.
Wir entwickeln die Theorie rigoros: Jensens Formel, Kroneckers Satz ($\M(\alpha)=1$
genau dann, wenn $\alpha$ Null oder Einheitswurzel), Salem-Zahlen,
Dobrowolskis untere Schranke $\M(\alpha) \geq 1 + c(\log\log d/\log d)^3$,
die Salem-Boyd-Vermutung und die Verbindung zum Entropie-Begriff algebraischer
dynamischer Systeme (Lind-Schmidt-Ward).
\end{abstract}

\maketitle

\tableofcontents

% ==============================================================================
\section{Einleitung}
% ==============================================================================

In seiner Arbeit von 1933 zur Faktorisierung großer Zahlen \cite{Lehmer1933} führte
D.\,H.\ Lehmer das Polynom
\[
  L(x) = x^{10}+x^9-x^7-x^6-x^5-x^4-x^3+x+1
\]
ein und bemerkte, dass es ein ungewöhnlich kleines Mahler-Maß besitzt:
$\M(L) \approx 1{,}17628\ldots$.
Er fragte:
\begin{quote}
  \textit{Gibt es ein Polynom mit ganzzahligen Koeffizienten, dessen Mahler-Maß
  strikt zwischen $1$ und $\M(L)$ liegt?}
\end{quote}

Diese Frage, heute als \emph{Lehmers Problem} bekannt, ist eines der ältesten ungelösten
Probleme der algebraischen Zahlentheorie. Sie hat tiefe Verbindungen zu:
\begin{itemize}
  \item der Verteilung algebraischer ganzer Zahlen auf und nahe dem Einheitskreis;
  \item Salem-Zahlen und Pisot-Vijayaraghavan-Zahlen (PV-Zahlen);
  \item der Entropie algebraischer $\Z^d$-Wirkungen;
  \item der Weil-Höhe algebraischer Zahlen;
  \item der Spektraltheorie von Graphen.
\end{itemize}

\subsection{Aufbau des Aufsatzes}
Abschnitt~\ref{sec:mahler_def} definiert das Mahler-Maß und beweist Jensens Formel.
Abschnitt~\ref{sec:kronecker} beweist Kroneckers Satz.
Abschnitt~\ref{sec:lehmer_poly} analysiert Lehmers Polynom und Salem-Zahlen.
Abschnitt~\ref{sec:vermutung} formuliert die Lehmer-Vermutung präzise.
Abschnitt~\ref{sec:dobrowolski} beweist Dobrowolskis untere Schranke.
Abschnitt~\ref{sec:salem_boyd} behandelt die Salem-Boyd-Vermutung.
Abschnitt~\ref{sec:entropie} verbindet das Mahler-Maß mit dynamischer Entropie.
Abschnitt~\ref{sec:pisot} bespricht Pisot- und Salem-Zahlen.
Abschnitt~\ref{sec:mehrere_variablen} behandelt das mehrdimensionale Mahler-Maß.
Abschnitt~\ref{sec:berechnungen} fasst rechnerische Ergebnisse zusammen.
Abschnitt~\ref{sec:offene_probleme} listet offene Probleme auf.
Abschnitt~\ref{sec:fazit} schließt den Aufsatz ab.

% ==============================================================================
\section{Mahler-Maß: Definition und Grundeigenschaften}\label{sec:mahler_def}
% ==============================================================================

\begin{definition_de}[Mahler-Maß eines Polynoms]
  Sei $P(x) = a_d x^d + \cdots + a_0 \in \C[x]$ mit $a_d \neq 0$ und Nullstellen
  $\alpha_1,\ldots,\alpha_d \in \C$. Das \emph{Mahler-Maß} von $P$ ist
  \[
    \M(P) \;:=\; |a_d| \prod_{j=1}^d \max(1, |\alpha_j|).
  \]
  Das \emph{logarithmische Mahler-Maß} ist $\m(P) = \log \M(P)$.
\end{definition_de}

\begin{satz}[Jensensche Formel]
  Für $P(x) = a_d \prod_{j=1}^d (x - \alpha_j) \in \C[x]$ gilt:
  \[
    \M(P) \;=\; \exp\!\left(\int_0^1 \log|P(e^{2\pi i t})|\,dt\right).
  \]
\end{satz}

\begin{proof}
  Nach der Jensen-Formel für jeden Linearfaktor $(x - \alpha_j)$ gilt:
  \[
    \int_0^1 \log|e^{2\pi it} - \alpha_j|\,dt = \log\max(1, |\alpha_j|).
  \]
  Einsetzen in die Produktformel ergibt die Behauptung.
\end{proof}

\begin{proposition}[Grundeigenschaften des Mahler-Maßes]
  \begin{enumerate}[label=(\roman*)]
    \item $\M(PQ) = \M(P)\M(Q)$ (Multiplikativität);
    \item $\M(P) = \M(P^\ast)$, wo $P^\ast(x) = x^d \overline{P(1/\bar{x})}$
      das reziproke Polynom ist;
    \item $\M(P) \geq 1$ für monisches $P \in \Z[x]$ mit $P(0) \neq 0$.
  \end{enumerate}
\end{proposition}

\begin{definition_de}[Mahler-Maß einer algebraischen Zahl]
  Für eine algebraische ganze Zahl $\alpha$ mit monischem irreduziblem
  Minimalpolynom $P_\alpha \in \Z[x]$ sei $\M(\alpha) := \M(P_\alpha)$.
\end{definition_de}

% ==============================================================================
\section{Kroneckers Satz}\label{sec:kronecker}
% ==============================================================================

\begin{satz}[Kronecker, 1857 \cite{Kronecker1857}]\label{satz:kronecker}
  Sei $\alpha$ eine von Null verschiedene algebraische ganze Zahl.
  Dann gilt $\M(\alpha) = 1$ genau dann, wenn $\alpha$ eine Einheitswurzel ist.
  Äquivalent: $\M(P) = 1$ für ein monisches $P \in \Z[x]$ mit $P(0) \neq 0$ genau dann,
  wenn alle Nullstellen von $P$ auf dem Einheitskreis liegen.
\end{satz}

\begin{proof}
  \emph{($\Leftarrow$)}: Liegen alle Nullstellen auf dem Einheitskreis, so gilt
  $\max(1,|\alpha_j|) = 1$ für alle $j$, also $\M(P) = 1$.

  \emph{($\Rightarrow$)}: Gilt $\M(P) = 1$ und ist $P$ monisch in $\Z[x]$ mit
  $|\alpha_j| \leq 1$ für alle $j$, so folgt $|a_0| = |\prod \alpha_j| \leq 1$;
  da $a_0 \in \Z \setminus \{0\}$, muss $|a_0| = 1$, also alle $|\alpha_j| = 1$.
  Nun erzeugen die $n$-ten Potenzen $\alpha_j^n$ Polynome mit beschränkten
  ganzzahligen Koeffizienten; von Kroneckers Endlichkeitsargument folgt,
  dass $\alpha_j^m = \alpha_j^n$ für $m \neq n$, also sind alle $\alpha_j$
  Einheitswurzeln.
\end{proof}

\begin{korollar}
  Für ein Kreisteilungspolynom $\Phi_n(x)$ gilt $\M(\Phi_n) = 1$.
\end{korollar}

% ==============================================================================
\section{Lehmers Polynom und Salem-Zahlen}\label{sec:lehmer_poly}
% ==============================================================================

\begin{definition_de}[Salem-Zahl]
  Eine \emph{Salem-Zahl} ist eine reelle algebraische ganze Zahl $\tau > 1$, deren
  Konjugierte alle auf dem oder im Innern des Einheitskreises liegen, wobei mindestens
  eine auf dem Einheitskreis liegt.
\end{definition_de}

\begin{definition_de}[Pisot-Vijayaraghavan-Zahl]
  Eine \emph{PV-Zahl} ist eine reelle algebraische ganze Zahl $\theta > 1$, deren
  sämtliche Konjugierte den Betrag $< 1$ haben.
\end{definition_de}

\begin{example}[Lehmers Polynom]
  Das Polynom
  \[
    L(x) = x^{10}+x^9-x^7-x^6-x^5-x^4-x^3+x+1
  \]
  ist irreduzibel über $\Q$ und besitzt:
  \begin{itemize}
    \item eine reelle Nullstelle $\lambda_{10} > 1$ (eine Salem-Zahl);
    \item eine reelle Nullstelle $\lambda_{10}^{-1} \in (0,1)$;
    \item $8$ komplexe Nullstellen (nahe dem) auf dem Einheitskreis, in $4$
      konjugierten Paaren.
  \end{itemize}
  Das Mahler-Maß beträgt $\M(L) = \lambda_{10} \approx 1{,}1762808182599175\ldots$.
\end{example}

\begin{bemerkung}
  Das Polynom $L(x)$ ist reziprok: $x^{10}L(1/x) = L(x)$.
  Dies ist charakteristisch für Salem-Zahlen.
\end{bemerkung}

% ==============================================================================
\section{Die Lehmer-Vermutung}\label{sec:vermutung}
% ==============================================================================

\begin{vermutung}[Lehmer-Vermutung \cite{Lehmer1933}]\label{verm:lehmer}
  Es gibt eine absolute Konstante $c > 1$, so dass für jede algebraische ganze Zahl
  $\alpha$, die keine Einheitswurzel ist, gilt:
  \[
    \M(\alpha) \;\geq\; c.
  \]
  Stärkere Form: Für jedes monische $P \in \Z[x]$ mit $P(0) \neq 0$, das kein Produkt
  von Kreisteilungspolynomen ist, gilt:
  \[
    \M(P) \;\geq\; \lambda_{10} \approx 1{,}17628.
  \]
\end{vermutung}

\begin{bemerkung}
  Die \emph{schwache} Form der Vermutung besagt nur, dass $1$ kein Häufungspunkt der
  Menge $\{\M(\alpha) : \alpha \text{ alg.\ ganzzahlig, keine Einheitswurzel}\}$ ist.
  Die \emph{starke} Form behauptet, dass $\lambda_{10}$ das Minimum ist.
\end{bemerkung}

\begin{satz}[Smyth \cite{Smyth1971}]
  Für ein nicht-reziprokes monisches Polynom $P \in \Z[x]$ (d.\,h.\ $P \neq x^{\deg P}P(1/x)$),
  das kein Produkt von Kreisteilungspolynomen ist, gilt:
  \[
    \M(P) \;\geq\; \theta_3 \;=\; 1{,}3247179572\ldots,
  \]
  wobei $\theta_3$ die kleinste PV-Zahl (reelle Nullstelle von $x^3-x-1$) ist.
\end{satz}

% ==============================================================================
\section{Dobrowolskis untere Schranke}\label{sec:dobrowolski}
% ==============================================================================

\begin{satz}[Dobrowolski \cite{Dobrowolski1979}]\label{satz:dobro}
  Sei $\alpha$ eine algebraische ganze Zahl vom Grad $d \geq 2$ über $\Q$, die keine
  Einheitswurzel ist. Dann gilt:
  \[
    \M(\alpha) \;\geq\; 1 + \frac{1}{1200}\left(\frac{\log\log d}{\log d}\right)^3.
  \]
\end{satz}

\begin{bemerkung}
  Verbesserte Konstanten wurden von Voutier \cite{Voutier1996} und Cantor-Straus
  \cite{CantorStraus1982} erzielt. Die beste bekannte Form ist
  $\M(\alpha) \geq 1 + (1-\varepsilon(d))(\log\log d / \log d)^3$
  mit $\varepsilon(d) \to 0$.
\end{bemerkung}

\textbf{Schlüsselideen des Beweises:}

\begin{lemma}[Hilfspolynomabschätzung]
  Für eine Primzahl $q$ und eine algebraische ganze Zahl $\alpha$ vom Grad $d$ betrachte
  \[
    P_q(x) = \prod_{\zeta^q = 1} P_\alpha(\zeta x) \in \Z[x].
  \]
  Es gilt $\log \M(P_q) \leq q \cdot \log \M(\alpha)$,
  während $P_q$ eine große $\ell^\infty$-Norm hat, falls $\M(\alpha) < 2$.
\end{lemma}

\begin{proof}[Beweisskizze von Satz~\ref{satz:dobro}]
  Nach dem Bertrandschen Postulat gibt es eine Primzahl $q \in [\frac{\log d}{2}, \log d]$.
  Die Frobenius-Wirkung liefert ein Polynom der Norm $\geq q$, während das
  Mahler-Maß von $P_q$ höchstens $\M(\alpha)^q$ ist. Der Widerspruch bei zu kleinem
  $\M(\alpha)$ ergibt die untere Schranke.
\end{proof}

% ==============================================================================
\section{Die Salem-Boyd-Vermutung}\label{sec:salem_boyd}
% ==============================================================================

\begin{vermutung}[Salem-Boyd \cite{Boyd1977}]\label{verm:salemboyd}
  Die kleinste Salem-Zahl ist die Lehmer-Zahl $\lambda_{10} \approx 1{,}17628$.
  Äquivalent: Lehmers Polynom $L(x)$ liefert das kleinste Mahler-Maß $> 1$ unter allen
  monischen ganzzahligen Polynomen.
\end{vermutung}

\begin{bemerkung}
  Boyd \cite{Boyd1977,Boyd1981} hat alle Salem-Zahlen vom Grad $\leq 40$ und
  Mahler-Maß $\leq 1{,}3$ berechnet und dabei stets $\lambda_{10}$ als das kleinste
  bestätigt. Die zweitkleinste bekannte Salem-Zahl hat
  $\M \approx 1{,}1884\ldots$ (Grad 18).
\end{bemerkung}

\begin{satz}[Mossinghoff-Rhin-Wu \cite{MRW2008}]
  Die Lehmer-Vermutung gilt für alle Polynome vom Grad $\leq 44$:
  Ist $P \in \Z[x]$, $\deg P \leq 44$, $P(0) \neq 0$ und $P$ kein Produkt von
  Kreisteilungspolynomen, so gilt $\M(P) \geq \lambda_{10}$.
\end{satz}

% ==============================================================================
\section{Entropie algebraischer dynamischer Systeme}\label{sec:entropie}
% ==============================================================================

\begin{definition_de}[Algebraische $\Z^d$-Wirkung]
  Für ein Laurent-Polynom $f \in \Z[x_1^{\pm 1},\ldots,x_d^{\pm 1}]$ ist
  $X_f := \widehat{\Z[x_1^{\pm 1},\ldots,x_d^{\pm 1}]/(f)}$
  die Pontrjagin-duale kompakte abelsche Gruppe.
  Die $\Z^d$-Wirkung auf $X_f$ ist durch duale Verschiebungen gegeben.
\end{definition_de}

\begin{satz}[Lind-Schmidt-Ward \cite{LSW1990}]\label{satz:lsw}
  Die topologische Entropie der $\Z^d$-Wirkung auf $X_f$ ist gleich dem
  (mehrdimensionalen) logarithmischen Mahler-Maß:
  \[
    h(X_f) \;=\; \m(f) \;=\; \int_{\T^d} \log|f(e^{2\pi i \theta_1},\ldots,e^{2\pi i \theta_d})|\,
    d\theta_1 \cdots d\theta_d.
  \]
\end{satz}

\begin{bemerkung}
  Satz~\ref{satz:lsw} bedeutet: Lehmers Problem ist äquivalent zur Frage, ob es eine
  algebraische $\Z$-Wirkung mit topologischer Entropie in $(0, \log\lambda_{10})$ gibt.
  Die Antwort ist vermutlich \emph{nein}.
\end{bemerkung}

\begin{example}[Smyths Formel]
  Für $f = 1 + x + y \in \Z[x^{\pm 1}, y^{\pm 1}]$ gilt
  \[
    \m(f) \;=\; \frac{3\sqrt{3}}{4\pi} L(\chi_{-3}, 2)
           \;=\; L'(\chi_{-3}, -1),
  \]
  wobei $\chi_{-3}$ der primitive Dirichlet-Charakter vom Führer $3$ ist.
\end{example}

% ==============================================================================
\section{Pisot-Zahlen, Salem-Zahlen und Häufungspunkte}\label{sec:pisot}
% ==============================================================================

\begin{satz}[Siegel 1944 \cite{Siegel1944}]
  Die kleinste PV-Zahl ist die reelle Nullstelle von $x^3 - x - 1$
  (die \emph{Kunststoffzahl} $\theta_3 \approx 1{,}3247$).
\end{satz}

\begin{satz}[Salem 1944 \cite{Salem1944}]
  Jede PV-Zahl größer als $1$ ist Häufungspunkt (von oben und unten) von Salem-Zahlen.
  Umgekehrt ist jeder Häufungspunkt von Salem-Zahlen eine PV-Zahl.
\end{satz}

\begin{proposition}[Charakterisierung durch Minimalpolynome]
  \begin{enumerate}[label=(\roman*)]
    \item Das Minimalpolynom einer Salem-Zahl ist reziprok (d.\,h.\ $x^d P(1/x) = P(x)$)
      und hat geraden Grad.
    \item Das Minimalpolynom einer PV-Zahl ist nicht reziprok.
    \item Lehmers Polynom $L(x)$ ist reziprok: $x^{10}L(1/x) = L(x)$.
  \end{enumerate}
\end{proposition}

% ==============================================================================
\section{Mehrdimensionales Mahler-Maß}\label{sec:mehrere_variablen}
% ==============================================================================

\begin{definition_de}[Mehrdimensionales Mahler-Maß]
  Für $f \in \C[x_1^{\pm 1},\ldots,x_n^{\pm 1}]$ sei
  \[
    \m(f) \;:=\; \int_0^1\cdots\int_0^1
    \log|f(e^{2\pi it_1},\ldots,e^{2\pi it_n})|\,dt_1\cdots dt_n.
  \]
\end{definition_de}

\begin{example}[Deningers Formel \cite{Deninger1997}]
  \[
    \m(1+x+1/x+y+1/y) \;=\; \frac{15}{(2\pi)^2}\, L(E_{15}, 2),
  \]
  wobei $E_{15}$ die elliptische Kurve $y^2 + xy + y = x^3 + x^2 - 10x - 10$
  vom Führer $15$ ist.
\end{example}

\begin{bemerkung}
  Deninger \cite{Deninger1997} vermutete eine allgemeine Verbindung: Für viele
  zweidimensionale Polynome $f$, die eine elliptische Kurve $E_f$ definieren, ist
  $\m(f)$ ein rationales Vielfaches von $L'(E_f, 0)$.
  Dies wurde für Hunderte von Beispielen numerisch bestätigt (Rodriguez Villegas
  \cite{RV1999}) und in einigen Fällen bewiesen (Rogers-Zudilin \cite{RZ2012}).
\end{bemerkung}

% ==============================================================================
\section{Rechnerische Ergebnisse}\label{sec:berechnungen}
% ==============================================================================

\begin{satz}[Mossinghoff \cite{Mossinghoff1998}]
  Die folgende Tabelle zeigt die kleinsten bekannten Mahler-Maße $> 1$
  für Polynome in $\Z[x]$:
  \begin{center}
    \begin{tabular}{lll}
      \toprule
      Rang & $\M$ & Polynom \\
      \midrule
      1 & $1{,}17628\ldots$ & $x^{10}+x^9-x^7-\cdots+1$ (Lehmer) \\
      2 & $1{,}18836\ldots$ & Salem-Zahl vom Grad 18 \\
      3 & $1{,}20002\ldots$ & Salem-Zahl vom Grad 20 \\
      \bottomrule
    \end{tabular}
  \end{center}
\end{satz}

% ==============================================================================
\section{Weil-Höhe und Lehmers Problem}\label{sec:weil_hoehe}
% ==============================================================================

\begin{definition_de}[Weil-Höhe]
  Für eine algebraische Zahl $\alpha$ mit Minimalpolynom $P_\alpha \in \Z[x]$ (monisch,
  irreduzibel) ist die \emph{logarithmische Weil-Höhe}:
  \[
    h(\alpha) \;:=\; \frac{1}{\deg\alpha} \log \M(P_\alpha)
              \;=\; \frac{\m(P_\alpha)}{\deg\alpha}.
  \]
\end{definition_de}

\begin{vermutung}[Lehmer-Vermutung über Weil-Höhe]\label{verm:lehmer_hoehe}
  Es gibt eine absolute Konstante $c > 0$, so dass für jede algebraische Zahl
  $\alpha \neq 0$, die keine Einheitswurzel ist, gilt:
  \[
    h(\alpha) \;\geq\; \frac{c}{\deg\alpha}.
  \]
\end{vermutung}

\begin{bemerkung}
  Diese Formulierung von Lehmers Problem (mit der Weil-Höhe) ist zur ursprünglichen
  äquivalent. Sie ist besonders nützlich in der Geometrie der Zahlen und der
  Arithmetischen Geometrie.
\end{bemerkung}

% ==============================================================================
\section{Offene Probleme}\label{sec:offene_probleme}
% ==============================================================================

\begin{vermutung}[Starke Lehmer-Vermutung \cite{Lehmer1933}]
  Für jedes monische $P \in \Z[x]$ mit $P(0) \neq 0$, das kein Produkt von
  Kreisteilungspolynomen ist, gilt $\M(P) \geq \lambda_{10} = 1{,}1762808\ldots$.
\end{vermutung}

\begin{vermutung}[Mehrdimensionale Lehmer-Vermutung \cite{Boyd1981}]
  Das Mahler-Maß eines mehrfachen Polynoms in $\Z[x_1^\pm,\ldots,x_n^\pm]$,
  das auf dem Torus $\T^n$ nicht verschwindet, ist mindestens $\lambda_{10}$.
\end{vermutung}

\begin{vermutung}[Bogomolov-Eigenschaft]
  Für jeden Zahlkörper $K$ gibt es $c(K) > 0$, so dass für alle
  $\alpha \in K^\times$, die keine Einheitswurzeln sind:
  $h(\alpha) \geq c(K)/[K:\Q]$.
\end{vermutung}

\begin{bemerkung}
  Die Bogomolov-Eigenschaft für $K = \Q$ ist äquivalent zur Lehmer-Vermutung.
  Sie ist bekannt für abelsche Erweiterungen von $\Q$
  (Amoroso-Dvornicich \cite{AD2000}).
\end{bemerkung}

% ==============================================================================
\section{Fazit}\label{sec:fazit}
% ==============================================================================

Lehmers Problem verbindet, trotz seiner elementaren Formulierung, tiefe Aspekte der
algebraischen Zahlentheorie: die Geometrie algebraischer ganzer Zahlen am Einheitskreis,
die Arithmetik von Salem- und Pisot-Zahlen sowie die Entropie algebraischer dynamischer
Systeme. Die einzige bewiesene untere Schranke von Substanz ist Dobrowolskis Ergebnis,
das weit hinter dem vermuteten Schwellenwert $\lambda_{10}$ zurückbleibt.
Lehmers Polynom hält nach 90 Jahren den Rekord, und die Vermutung, dass kein Mahler-Maß
strikt zwischen $1$ und $\lambda_{10}$ liegt, gilt als eines der schwierigsten offenen
Probleme der algebraischen Zahlentheorie. Die unerwarteten Verbindungen zum
mehrdimensionalen Mahler-Maß, zu $L$-Funktionen und elliptischen Kurven deuten auf eine
tiefe arithmetische Struktur hin.

\begin{thebibliography}{99}

\bibitem{AD2000}
F.\ Amoroso, R.\ Dvornicich,
\textit{A lower bound for the height in abelian extensions},
J.\ Number Theory \textbf{80} (2000), Nr.\ 2, 260--272.

\bibitem{Boyd1977}
D.\,W.\ Boyd,
\textit{Small Salem numbers},
Duke Math.\ J.\ \textbf{44} (1977), Nr.\ 2, 315--328.

\bibitem{Boyd1981}
D.\,W.\ Boyd,
\textit{Speculations concerning the range of Mahler's measure},
Canad.\ Math.\ Bull.\ \textbf{24} (1981), Nr.\ 4, 453--469.

\bibitem{CantorStraus1982}
D.\,G.\ Cantor, E.\,G.\ Straus,
\textit{On a conjecture of D.\,H.\ Lehmer},
Acta Arith.\ \textbf{42} (1982), Nr.\ 2, 97--100.

\bibitem{Deninger1997}
C.\ Deninger,
\textit{Deligne periods of mixed motives, $K$-theory and the entropy of certain
$\mathbb{Z}^n$-actions},
J.\ Amer.\ Math.\ Soc.\ \textbf{10} (1997), Nr.\ 2, 259--281.

\bibitem{Dobrowolski1979}
E.\ Dobrowolski,
\textit{On a question of Lehmer and the number of irreducible factors of a polynomial},
Acta Arith.\ \textbf{34} (1979), Nr.\ 4, 391--401.

\bibitem{Kronecker1857}
L.\ Kronecker,
\textit{Zwei Sätze über Gleichungen mit ganzzahligen Koeffizienten},
J.\ Reine Angew.\ Math.\ \textbf{53} (1857), 173--175.

\bibitem{Lehmer1933}
D.\,H.\ Lehmer,
\textit{Factorization of certain cyclotomic functions},
Ann.\ of Math.\ (2) \textbf{34} (1933), Nr.\ 3, 461--479.

\bibitem{LSW1990}
D.\,A.\ Lind, K.\ Schmidt, T.\ Ward,
\textit{Mahler measure and entropy for commuting automorphisms of compact groups},
Invent.\ Math.\ \textbf{101} (1990), Nr.\ 3, 593--629.

\bibitem{Mossinghoff1998}
M.\,J.\ Mossinghoff,
\textit{Polynomials with small Mahler measure},
Math.\ Comp.\ \textbf{67} (1998), Nr.\ 224, 1697--1705.

\bibitem{MRW2008}
M.\,J.\ Mossinghoff, G.\ Rhin, Q.\ Wu,
\textit{Minimal Mahler measures},
Experiment.\ Math.\ \textbf{17} (2008), Nr.\ 4, 451--458.

\bibitem{RV1999}
F.\ Rodriguez Villegas,
\textit{Modular Mahler measures I},
Math.\ Appl.\ \textbf{467}, Kluwer (1999), 17--48.

\bibitem{RZ2012}
M.\ Rogers, W.\ Zudilin,
\textit{From $L$-series of elliptic curves to Mahler measures},
Compos.\ Math.\ \textbf{148} (2012), Nr.\ 2, 385--414.

\bibitem{Salem1944}
R.\ Salem,
\textit{A remarkable class of algebraic integers},
Duke Math.\ J.\ \textbf{11} (1944), 103--108.

\bibitem{Siegel1944}
C.\,L.\ Siegel,
\textit{Algebraic integers whose conjugates lie in the unit circle},
Duke Math.\ J.\ \textbf{11} (1944), 597--602.

\bibitem{Smyth1971}
C.\,J.\ Smyth,
\textit{On the product of the conjugates outside the unit circle of an algebraic integer},
Bull.\ London Math.\ Soc.\ \textbf{3} (1971), 169--175.

\bibitem{Smyth1981}
C.\,J.\ Smyth,
\textit{On measures of polynomials in several variables},
Bull.\ Austral.\ Math.\ Soc.\ \textbf{23} (1981), Nr.\ 1, 49--63.

\bibitem{Voutier1996}
P.\ Voutier,
\textit{An effective lower bound for the height of algebraic numbers},
Acta Arith.\ \textbf{74} (1996), Nr.\ 1, 81--95.

\end{thebibliography}

\end{document}
